\documentclass[11pt,a4paper]{article}

% ============================================================
% PACOTES
% ============================================================
\usepackage[utf8]{inputenc}
\usepackage[T1]{fontenc}
\usepackage[brazil]{babel}

\usepackage{amsmath,amssymb,amsfonts,mathtools}
\usepackage{bm}
\usepackage{physics}
\usepackage{siunitx}
\usepackage{booktabs}
\usepackage{geometry}
\usepackage{float}
\usepackage{tikz}
\usepackage{pgfplots}
\usepackage{listings}
\usepackage{xcolor}
\usepackage{hyperref}
\usetikzlibrary{intersections}

\pgfplotsset{compat=1.18}
\geometry{margin=2.5cm}

\hypersetup{
    colorlinks=true,
    linkcolor=blue,
    citecolor=blue,
    urlcolor=blue
}

% ============================================================
% COMANDOS
% ============================================================
\newcommand{\pdvT}[2]{\left(\frac{\partial #1}{\partial #2}\right)_T}
\newcommand{\pdvV}[2]{\left(\frac{\partial #1}{\partial #2}\right)_v}
\newcommand{\pdvS}[2]{\left(\frac{\partial #1}{\partial #2}\right)_s}
\newcommand{\pd}[2]{\frac{\partial #1}{\partial #2}}
\newcommand{\pdd}[2]{\frac{\partial^2 #1}{\partial #2^2}}
\newcommand{\e}[1]{\mathrm{e}^{#1}}
\renewcommand{\dd}{\mathrm{d}}
\newcommand{\tildep}{\widetilde{p}}
\newcommand{\tildev}{\widetilde{v}}
\newcommand{\tildeT}{\widetilde{T}}

% ============================================================
% DOCUMENTO
% ============================================================
\begin{document}
%
\begin{center} 
    {\Large \textbf{Primeira série de exercícios}}\\[0.6cm]
    {\large Tópicos de Mecânica Estatística -- Transições de Fases -- 2026}\\[0.2cm]
    {\large Fluido de van der Waals -- 12/8/2026}\\[0.4cm]    
    {\large Data de Entrega -- 25/8/2026}\\[0.2cm]
    {\large Data limite -- 28/8/2026}\\[0.4cm]
    {\Large \textbf{Autor:} Thiago Siqueira Domingues (NºUSP: 8942194)}    
\end{center}
%
\vspace{0.5cm}
% ============================================================
\section*{Problema (1) -- Equação de estados correspondentes}
%
A equação de estado de van der Waals é
%
\begin{equation}
    \left(p+\frac{a}{v^2}\right)(v-b)=RT
\label{eq:vdw}
\end{equation}
%
ou, equivalentemente,
%
\begin{equation}
    p(T,v)=\frac{RT}{v-b}-\frac{a}{v^2}.
\label{eq:pvdw}
\end{equation}
%
Aqui, $p$ é a pressão, $v$ é o volume molar, $T$ é a temperatura absoluta e $R$ é a constante universal dos gases. Os parâmetros fenomenológicos $a>0$ e $b>0$ representam, respectivamente, a atração média entre as moléculas e o volume excluído pelas moléculas.
%
Definindo as variáveis adimensionais, 
%
\begin{equation}
    \tilde p=\frac{p}{p_c},
    \qquad
    \tilde v=\frac{v}{v_c},
    \qquad
    \tilde T=\frac{T}{T_c},
\label{eq:reduced_variables}
\end{equation}
%
em que $p_c$, $v_c$ e $T_c$ são os valores da pressão, do volume por mol ou molar e da temperatura no ponto crítico. 
%
Agora, devemos obter a \textit{equação de ``estados correspondentes"}, válida para qualquer fluído. Primeiro, devemos calcular os parâmetros críticos $p_c$, $v_c$ e $T_c$ em função dos parâmetros fenomenológicos $a$ e $b$. Existem duas maneiras de fazer isto, até onde eu conheça. Seguindo a maneira apresentada em aula e que também pode ser vista em~\cite{salinas}, podemos reescrever a equação de van der Waals na forma de um polinômio cúbico em $v$. 
%
Então, multiplicando a Eq.~\ref{eq:vdw} for $v^2$, temos  
%
\begin{equation*}
    (pv^2 + a) (v-b) = v^2 RT \Rightarrow pv^3 - pv^2b + av - ab = v^2RT
\end{equation*} 
%
Agora, dividindo por $p$, chegamos em
%
\begin{equation*}
    v^3 - v^2b + \frac{av}{p} - \frac{ab}{p} - \frac{v^2 RT}{p} = 0,
\end{equation*}
%
agrupando os termos,
%
\begin{equation}
    v^3-\left(b+\frac{RT_c}{p_c}\right)v^2
    +\frac{a}{p_c}v
    -\frac{ab}{p_c}=0.
\label{eq:cubic_polynom}
\end{equation}
%
Agora, no ponto crítico ($T = T_c$, $p = p_c$) as três raízes deste polinômio são idênticas, ou seja,
%
\begin{equation*}
    v_1=v_2=v_3=v_c.
\end{equation*}
%
Portanto, o polinômio cúbico pode ser escrito, no ponto crítico, como
%
\begin{equation}
    (v-v_c)^3=0,
\end{equation}
%
ou, explicitamente,
%
\begin{equation}
    v^3-3v_c v^2+3v_c^2v-v_c^3=0.
\end{equation}
%
Igualando esta equação com a Eq.~\ref{eq:cubic_polynom} e comparando os coeficientes dos dois polinômios, obtemos o sistema
%
\begin{align}
    3v_c &= b+\frac{RT_c}{p_c},
    \label{eq:critical_comparison_1}\\
    3v_c^2 &= \frac{a}{p_c},
    \label{eq:critical_comparison_2}\\
    v_c^3 &= \frac{ab}{p_c}.
    \label{eq:critical_comparison_3}
\end{align}
%
As duas últimas equações permitem determinar diretamente $v_c$ e $p_c$. 
Da Eq.~\ref{eq:critical_comparison_2}, temos
%
\begin{equation}
    p_c=\frac{a}{3v_c^2}.
\end{equation}
%
Substituindo esse resultado na Eq.~\ref{eq:critical_comparison_3}, encontramos
%
\begin{equation}
    v_c^3
    =
    \frac{ab}{a/(3v_c^2)}
    =
    3bv_c^2.
\end{equation}
%
Como $v_c\neq 0$, segue que
%
\begin{equation}
    \boxed{v_c=3b}.
\label{eq:vc}
\end{equation}
%
Substituindo esse resultado em
%
\begin{equation}
    p_c=\frac{a}{3v_c^2},
\end{equation}
%
obtemos
%
\begin{equation}
    \boxed{p_c=\frac{a}{27b^2}}.
\label{eq:pc}
\end{equation}
%
Finalmente, podemos determinar a temperatura crítica utilizando a
Eq.~\ref{eq:critical_comparison_1}. Temos
%
\begin{equation}
    3v_c=b+\frac{RT_c}{p_c},
\end{equation}
%
e, portanto,
%
\begin{equation}
    RT_c=p_c(3v_c-b).
\end{equation}
%
Usando $v_c=3b$ e $p_c=a/(27b^2)$, resulta
%
\begin{equation}
    RT_c
    =
    \frac{a}{27b^2}(9b-b)
    =
    \frac{8a}{27b}.
\end{equation}
%
Assim,
%
\begin{equation}
    \boxed{T_c=\frac{8a}{27Rb}}.
\label{eq:Tc}
\end{equation}
%
Portanto, os parâmetros críticos da equação de van der Waals são
%
\begin{equation}
    \boxed{
    v_c=3b,
    \qquad
    p_c=\frac{a}{27b^2},
    \qquad
    T_c=\frac{8a}{27Rb}.
    }
\label{eq:vdw_critical_parameters}
\end{equation}
%
Podemos agora substituir esses resultados na equação de estado
%
\begin{equation}
    p=\frac{RT}{v-b}-\frac{a}{v^2},
\end{equation}
%
e expressá-la exclusivamente em termos das variáveis reduzidas
%
\begin{equation}
    \tilde p=\frac{p}{p_c},
    \qquad
    \tilde v=\frac{v}{v_c},
    \qquad
    \tilde T=\frac{T}{T_c}.
\end{equation}
%
Como
%
\begin{equation}
    p=p_c\tilde p,
    \qquad
    v=v_c\tilde v=3b\tilde v,
    \qquad
    T=T_c\tilde T,
\end{equation}
%
a equação de van der Waals torna-se
%
\begin{equation}
    p_c\tilde p
    =
    \frac{RT_c\tilde T}{3b\tilde v-b}
    -
    \frac{a}{9b^2\tilde v^2}.
\end{equation}
%
Como
%
\begin{equation}
    RT_c=\frac{8a}{27b},
\end{equation}
%
temos
%
\begin{equation}
    p_c\tilde p
    =
    \frac{8a}{27b^2(3\tilde v-1)}\tilde T
    -
    \frac{a}{9b^2\tilde v^2}.
\end{equation}
%
Utilizando ainda
%
\begin{equation}
    p_c=\frac{a}{27b^2},
\end{equation}
%
podemos dividir toda a equação por $p_c$, obtendo
%
\begin{equation}
    \frac{a}{27b^2} \tilde p
    = \frac{8a}{27b^2(3 \tilde v - 1)} \tilde T - \frac{a}{9b^2 \tilde v^2} \Rightarrow  
    \tilde p = \frac{8\tilde T}{3\tilde v-1}
    -
    \frac{3}{\tilde v^2}.
\end{equation}
%
Agora, rearranjando os termos da equação ficamos com,
%
\begin{equation}
    \tilde p + \frac{3}{\tilde v^2}
    =
    \frac{8\tilde T}{3\tilde v-1} = \frac{8\tilde T}{3(\tilde v- \frac{1}{3})}     
\end{equation}
%
Multiplicando os dois lados por $(\tilde v - \frac{1}{3})$, finalmente obtemos a equação de estados correspondentes para um fluido de van der Waals:
%
\begin{equation}
    \boxed{
        \left(\tilde{p} + \frac{3}{\tilde{v}^2}\right) \left(\tilde{v} - \frac{1}{3}\right) = \frac{8}{3}\tilde{T}
    }.
    \label{eq:corresponding_states_vdw}
\end{equation}
%
Essa é a chamada \textbf{equação de estados correspondentes}. O resultado não depende de $a$ nem de $b$. Portanto, todos os fluidos descritos por essa aproximação possuem a mesma equação reduzida.
%
\begin{equation*}
    \boxed{\text{Universalidade: a equação reduzida não depende de }a,b.}
\end{equation*}
% ===========================================================
%
\section*{Problema (2) -- Energia livre de Helmholtz e a equação de van der Waals}
%
Partimos da primeira lei da termodinâmica para um sistema simples
compressível. Como estamos trabalhando com grandezas molares,
escrevemos a energia interna molar como uma função da entropia molar
e do volume molar,
%
\begin{equation}
    u=u(s,v).
\end{equation}
%
A primeira lei da termodinâmica, por mol, é
%
\begin{equation}
\dd u = T\,\dd s-p\,\dd v.
\label{eq:primeira_lei}
\end{equation}
%
Como $u=u(s,v)$, seu diferencial total pode ser escrito como
%
\begin{equation}
    \dd u
    =
    \left(\frac{\partial u}{\partial s}\right)_v
    \dd s
    +
    \left(\frac{\partial u}{\partial v}\right)_s
    \dd v.
\end{equation}
%
Comparando essa expressão com a primeira lei, Eq.~\eqref{eq:primeira_lei}, obtemos
%
\begin{equation}
    T=
    \left(\frac{\partial u}{\partial s}\right)_v
    \label{eq:T}
\end{equation}
%
e
%
\begin{equation}
    -p=
    \left(\frac{\partial u}{\partial v}\right)_s.
    \label{eq:p}
\end{equation}
%
Essas são as relações fundamentais que identificam $T$ e $-p$
como as variáveis conjugadas às variáveis extensivas específicas
$molares$ $s$ e $v$.  A energia interna possui como variáveis naturais $s$ e $v$,  $u=u(s,v)$. Entretanto, frequentemente é mais conveniente trabalhar com a temperatura $T$ em vez da entropia $s$.
Para substituir a variável $s$ por sua variável conjugada $T$,
fazemos uma transformação de Legendre e definimos a energia livre
de Helmholtz molar:
%
\begin{equation}
    f(T,v)=u-Ts.
    \label{eq:helmholtz}
\end{equation}
%
Equivalentemente, no equilíbrio termodinâmico, podemos escrever
a transformação de Legendre na forma:
%
\begin{equation}
    f(T,v)
    =
    \min_s
    \left\{
        u(s,v)-Ts
    \right\}.
    \label{eq:legendre}
\end{equation}
%
A partir da definição
%
\begin{equation}
    f=u-Ts,
\end{equation}
%
calculamos seu diferencial:
\begin{align}
    \dd f
    &=
    \dd u-\dd(Ts)\\
    &=
    \dd u-T\,\dd s-s\,\dd T.
\end{align}
%
Usando a primeira lei,
%
\begin{equation}
    \dd u=T\,\dd s-p\,\dd v,
\end{equation}
%
obtemos
%
\begin{align}
    \dd f
    &=
    T\,\dd s-p\,\dd v
    -T\,\dd s
    -s\,\dd T\\
    &=
    -s\,\dd T-p\,\dd v.
\end{align}
%
Portanto,
%
\begin{equation}
    \dd f=-s\,\dd T-p\,\dd v.
    \label{eq:df}
\end{equation}
%
Comparando com o diferencial total de uma função
$f=f(T,v)$,
%
\begin{equation}
    \dd f
    =
    \left(\frac{\partial f}{\partial T}\right)_v
    \dd T
    +
    \left(\frac{\partial f}{\partial v}\right)_T
    \dd v,
\end{equation}
%
obtemos as relações termodinâmicas
%
\begin{equation}
    \left(\frac{\partial f}{\partial T}\right)_v=-s
    \label{eq:s}
\end{equation}
%
e
%
\begin{equation}
    \left(\frac{\partial f}{\partial v}\right)_T=-p.
    \label{eq:pressure_free_energy}
\end{equation}
%
A segunda relação será utilizada para obter a energia livre de
Helmholtz a partir da equação de estado de van der Waals. Agora, substituindo a equação de estado de van der Waals para a pressão $p$:
%
\begin{equation}
    \left(\frac{\partial f}{\partial v}\right)_T = -\left( \frac{RT}{v-b} - \frac{a}{v^2} \right) = -\frac{RT}{v-b} + \frac{a}{v^2}
\end{equation}
%
Para encontrar $f(T,v)$, integramos ambos os lados  em relação ao volume molar $v$, mantendo a temperatura constante:
%
\begin{equation}
    \int
    \left(\frac{\partial f}{\partial v}\right)_T
    \dd v
    =
    -\int p\,\dd v.
\end{equation}
%
como a integração foi realizada mantendo $T$ constante, a ``constante'' de integração pode ser uma função de $T$. Chamaremos essa função de $f_0(T)$,
%
\begin{equation}
    f(T,v) = \int \left( -\frac{RT}{v-b} + \frac{a}{v^2} \right) dv +
    f_0(T),
\end{equation}
%
separamos a integral em duas partes:
%
\begin{equation}
    f(T,v) = -RT \int \frac{1}{v-b}\,dv + a \int v^{-2}\,dv +
    f_0(T)
\end{equation}
%
Resolvendo cada termo individualmente, para a primeira integral, utilizando a substituição 
%
$u = v-b$ (onde $du = dv$):
%
\begin{equation}
\int \frac{1}{v-b}\,dv = \ln(v-b)
\end{equation}
%
Para a segunda integral, aplicando a regra da potência:
%
\begin{equation}
\int v^{-2}\,dv = \frac{v^{-1}}{(-1)} = -\frac{1}{v}
\end{equation}
%
Reunindo os termos e adicionando a função de integração $f_0(T)$, que surge pelo fato de a integração ser em relação a $v$ com $T$ constante, chegamos em:
%
\begin{equation}
    \boxed{
    f = f(T,v)
    =
    -RT\ln(v-b)
    -\frac{a}{v}
    +f_0(T).
    }
    \label{eq:free_energy}
\end{equation}
%
Aqui, $f_0$ deve ser uma função bem comportada da temperatura $f_0 = f_0(T)$ e não pode ser determinada apenas a partir da equação de estado. De fato, $f_0=f_0(T)$ é determinada por propriedades térmicas adicionais, como o calor específico, como será visto no problema (8).
%
Para uma fase estável, a energia livre por mol $f$ deve ser convexa em relação às variáveis extensivas apropriadas. Em particular, em relação ao volume,
%
\begin{equation}
    \left(\pdd{f}{v}\right)_T\geq0.
\end{equation}
%
Essa condição pode ser relacionada diretamente à compressibilidade
isotérmica. Como
%
\begin{equation}
    \left(\pd{f}{v}\right)_T=-p,
\end{equation}
%
temos
%
\begin{equation}
    \left(\pdd{f}{v}\right)_T
    =
    -\left(\pd{p}{v}\right)_T.
\end{equation}
%
Assim, a estabilidade termodinâmica exige que
%
\begin{equation}
    \left(\pd{p}{v}\right)_T\leq0.
\end{equation}
%
Essa condição também segue da positividade da compressibilidade
isotérmica,
%
\begin{equation}
    \kappa_T
    =
    -\frac{1}{v}
    \left(\frac{\partial v}{\partial p}\right)_T
    >0.
\end{equation}
%
Para a equação de van der Waals,
%
\begin{equation}
    \left(\pd{p}{v}\right)_T
    =
    -\frac{RT}{(v-b)^2}
    +
    \frac{2a}{v^3}.
\end{equation}
%
Em determinadas temperaturas abaixo de $T_c$, essa quantidade torna-se
positiva em uma região intermediária de volumes. Consequentemente,
%
\begin{equation}
    \left(\pdd{f}{v}\right)_T<0
\end{equation}
%
nessa região. Substituindo a expressão da derivada da equação de van der Waals, a condição explícita para violar a desigualdade é:
%
\begin{equation}
    -\frac{RT}{(v-b)^2} + \frac{2a}{v^3} \geq 0
\end{equation}
%
Ou, isolando os termos de forma equivalente:
%
\begin{equation}
    \frac{RT}{(v-b)^2} \leq \frac{2a}{v^3}
\end{equation}
%
Logo, a energia livre por mol $f$ de van der Waals possui uma região não convexa (instável termodinamicamente). Esse é o problema da forma de energia livre notado por Maxwell.
%
\begin{equation*}
    \boxed{
    \text{A energia livre de van der Waals não é convexa em toda a região
    termodinamicamente acessível.}
    }
\end{equation*}
%
Essa região corresponde às famosas ``alças'' da isoterma de van der Waals~\cite{salinas}.
% ============================================================
\section*{Problema (3) -- Construção de Maxwell e dupla tangente}
%
Considere uma temperatura $T<T_c$. As fases líquida e gasosa coexistem nos volumes $v_L<v_G$. Na coexistência, ambas as fases possuem a mesma pressão,
%
\begin{equation}
    p_L=p_G=p_0,
\end{equation}
%
onde $p_0$ é a pressão de coexistência. A construção de Maxwell exige que
%
\begin{equation}    
    \int_{v_L}^{v_G}p(T,v)\,\dd v
    =
    p_0(v_G-v_L).
    \label{eq:maxwell}
\end{equation}
%
em que $v_G$ e $v_L$ são os volumes por mol das fases gasosa e líquida, que estão em coexistência a temperatura e pressão, respectivamente. Agora, utilizando a equação~\ref{eq:pressure_free_energy} podemos escrever que a integral de $p(T,v)$ em relação a $v$ é a própria energia livre de Helmholtz, 
%
\begin{equation}
   \int_{v_L}^{v_G} p(T,v)\, dv = - \int_{v_L}^{v_G} \left(\frac{\partial f}{\partial v}\right)_T\, dv
\end{equation}
%
Portanto, pelo teorema fundamental do cálculo, 
%
\begin{equation}
    \int_{v_L}^{v_G}p(T,v)\,\dd v
    =
    -[f_G-f_L],
\end{equation}
%
ou
%
\begin{equation}
    \int_{v_L}^{v_G}p(T,v)\,\dd v
    =
    f_L-f_G.
    \label{eq:intf}
\end{equation}
%
onde $f_L = f(v_L, T)$ e $f_G = f(v_G, T)$ são as energias livres das fases gasosa e líquida, respectivamente.
%
Usando a construção de Maxwell (Eq.~\ref{eq:maxwell}),
%
\begin{equation}
    f_L-f_G
    =
    p_0(v_G-v_L).
\end{equation}
%
Logo,
%
\begin{equation}
    \boxed{
    p_0
    =
    -\frac{f_G-f_L}{v_G-v_L}.
    }
    \label{eq:double_tangent}
\end{equation}
%
Essa equação da dupla tangente mostra que a pressão de coexistência
é determinada pela inclinação da reta que conecta os pontos
$(v_L,f_L)$ e $(v_G,f_G)$, veja Fig.~\ref{fig:double_tangent}. Essa é a dupla tangente proposta por Maxwell.
%
Agora, o gráfico de $f(T,v)$ contra v, abaixo da temperatura crítica ($T < T_c$), Fig.\ref{fig:double_tangent} podemos notar que a construção da dupla tangente restaura a convexidade da energia livre em relação ao volume. Além de passar pelos pontos $v_L$ e $v_G$, a reta deve ser tangente à energia livre nesses dois pontos. Na Fig.\ref{fig:equal_areas} vemos a correspondência desta construção no diagrama $p-v$.
%
\begin{figure}[H]
    \centering
    \includegraphics[width=0.7\linewidth]{Figures/statistical_mechanics/energia_livre.png}
    \caption{Gráfico da energia livre de Helmholtz por mol $f(T,v)$ em função do volume específico $v$, abaixo da temperatura crítica ($T < T_c$). Esta é uma representação esquemática da construção da dupla tangente. A envoltória convexa substitui a região não convexa da energia livre. A construção da dupla tangente é equivalente à construção de áreas iguais da figura~\ref{fig:equal_areas}. Figura baseada em~\cite{salinas}.}
    \label{fig:double_tangent}
\end{figure}
%
De maneira complementar, o diagrama $p-v$ representado na figura~\ref{fig:equal_areas} mostra que a construção da dupla tangente é equivalente a construção das áreas iguais de Maxwell.
%
\begin{figure}[htbp]
    \centering
    \includegraphics[width=0.7\linewidth]{Figures/statistical_mechanics/isoterma_vdW.png}
    \caption{Isoterma de van der Waals para $T<T_c$. A linha horizontal $p=p_0$ representa a pressão de coexistência. Os volumes $v_L$ e $v_G$ correspondem às fases líquida e gasosa. A construção de Maxwell exige que as áreas $A_1$ e $A_2$  na região hachurada sejam iguais. Figura baseada em~\cite{salinas}.
    }
\label{fig:equal_areas}
\end{figure}
% ============================================================
\section*{Problema (4) -- Expansão crítica: expoentes críticos de van der Waals}
%
Para obtermos os expoentes críticos de van der Waals, primeiro devemos defenir as seguintes quantidades (as chamadas variáveis reduzidas),
%
\begin{equation}
    \pi=\frac{p-p_c}{p_c},
    \qquad
    \omega=\frac{v-v_c}{v_c},
    \qquad
    t=\frac{T-T_c}{T_c}.
\end{equation}
%
Assim, $p=p_c(1+\pi)$, $v=v_c(1+\omega)=3b(1+\omega)$,
e $T=T_c(1+t)$. A equação reduzida obtida no problema 1 equação~\ref{eq:corresponding_states_vdw} pode ser reescrita na forma,
%
\begin{equation}
    \tilde p
    =
    \frac{8\tilde T}{3\tilde v-1}
    -
    \frac{3}{\tilde v^2},
\end{equation}
%
e como $\tilde p = 1 + \pi$, $\tilde v=1+\omega$ e $\tilde T = 1+t$,
%
obtemos $3\tilde v-1=2+3\omega$, $8 \tilde T = 8(1+t)$ e $\tilde v^2 = (1+ \omega)^2$.
%
Portanto, podemos reescrever a equação reduzida na forma,
%
\begin{equation}
    1+\pi
    =
    \frac{8(1+t)}{2+3\omega}
    -
    \frac{3}{(1+\omega)^2}.
\end{equation}
%
Logo,
%
\begin{equation}
    \pi
    =
    \frac{4(1+t)}{1+\frac32\omega}
    -
    \frac{3}{(1+\omega)^2}
    -1.
    \label{eq:pi_exact}
\end{equation}
%
Agora, para mostrar o comportamento nas proximidades do ponto crítico (quando $t, \omega \rightarrow 0$), devemos utilizar as expansões da Série de Maclaurin (série de Taylor centrada em x = 0)
%
\begin{equation}
    \frac{1}{1+x}
    =
    1-x+x^2-x^3+x^4+\mathcal{O}(x^5),
\end{equation}
%
\begin{equation}
    \frac{1}{(1+x)^2}
    =
    1-2x+3x^2-4x^3+5x^4+\mathcal{O}(x^5).
\end{equation}
%
Assim, a expansão próxima do ponto crítico é dada por,
%
\begin{equation}
    \frac{1}
    {1+\frac32\omega}
    =
    1
    -\frac32\omega
    +\frac94\omega^2
    -\frac{27}{8}\omega^3
    +\frac{81}{16}\omega^4
    +\mathcal O(\omega^5).
    \label{eq:first_expansion}
\end{equation}
%
então, 
%
\begin{align}
    \frac{4(1+t)}{1+\frac32\omega}
    &=
    4(1+t)
    \left[
    1-\frac32\omega
    +\frac94\omega^2
    -\frac{27}{8}\omega^3
    +\frac{81}{16}\omega^4
    +\mathcal{O}(\omega^5)
    \right],
\end{align}
%
consequentemente,
%
\begin{align}
    \frac{4(1+t)}{1+\frac32\omega}
    &=
    4+4t
    -6\omega
    -6t\omega
    +9\omega^2
    +9t\omega^2
    -\frac{27}{2}\omega^3
    -\frac{27}{2}t\omega^3
    +\frac{81}{4}\omega^4
    +\frac{81}{4}t\omega^4
    +\mathcal{O}(\omega^5).
\end{align}
%
enquanto
%
\begin{equation}
    (1+\omega)^{-2}
    =
    1-2\omega+3\omega^2
    -4\omega^3+5\omega^4
    +\mathcal O(\omega^5),
\end{equation}
%
portanto, 
%
\begin{equation}
    \frac{3}{(1+\omega)^2}
    =
    3-6\omega+9\omega^2-12\omega^3
    +15\omega^4
    +\mathcal{O}(\omega^5).
\end{equation}
%
Substituindo as expansões acima na Eq.~\ref{eq:pi_exact}, temos
%
\begin{align}
    \pi
    &=
    \left[
    4+4t
    -6\omega
    -6t\omega
    +9\omega^2
    +9t\omega^2
    -\frac{27}{2}\omega^3
    -\frac{27}{2}t\omega^3
    +\frac{81}{4}\omega^4
    +\frac{81}{4}t\omega^4
    \right]
    \nonumber\\
    &\quad
    -
    \left[
    3-6\omega+9\omega^2
    -12\omega^3+15\omega^4
    \right]
    -1
    +\mathcal{O}(\omega^5).
\end{align}
%
Organizando os termos de mesma ordem, obtemos
%
\begin{equation}
    \begin{aligned}
        \pi
        ={}&
        4t
        -6t\omega
        +9t\omega^2
        -\frac{27}{2}t\omega^3
        +\frac{81}{4}t\omega^4
        -\frac{3}{2}\omega^3
        +\frac{21}{4}\omega^4
        +\mathcal{O}(\omega^5).
    \end{aligned}
    \label{eq:pi_expanded}
\end{equation}
%
Como estamos interessados no comportamento nas proximidades do ponto
crítico, $t,\omega\rightarrow0$, podemos manter apenas os termos
dominantes. Nesse caso, na vizinhança crítica, a ordenação de escala relevante é
%
\begin{equation}
    \omega\sim |t|^{1/2}.
\end{equation}
%
Logo,
%
\begin{equation}
    t\omega^2\sim t^2,
\qquad
    t\omega^3\sim |t|^{5/2},
\qquad
    \omega^4\sim t^2.
\end{equation}
%
Consequentemente, ao nível solicitado,
%
\begin{equation}
    \boxed{
    \pi
    =
    4t
    -6t\omega
    -\frac32\omega^3
    +
    \mathcal{O}(\omega^4,t\omega^2).
    }
    \label{eq:critical_expansion}
\end{equation}
%
Essa é a expansão crítica fundamental do modelo de van der Waals.
% ============================================================
\section*{Problema (5) -- Forma assintótica da curva de coexistência de fases}
%
A partir da forma assintótica das isotermas de van der Waals para $t<0$ (abaixo do ponto crítico), na vizinhança imediata do ponto crítico obtida no exercício anterior (Eq.~\ref{eq:critical_expansion}), podemos escrever para as duas fases de coexistência,
%
\begin{align}
    \pi_1
    &=
    4t-6t\omega_1-\frac32\omega_1^3,
    \label{eq:pi1}
    \\
    \pi_2
    &=
    4t-6t\omega_2-\frac32\omega_2^3.
    \label{eq:pi2}
\end{align}
%
Na coexistência,

\begin{equation}
    \pi_1=\pi_2.
\end{equation}
%
Portanto,
%
\begin{equation}
    -6t(\omega_1-\omega_2)
    -
    \frac32(\omega_1^3-\omega_2^3)
    =0.
\end{equation}
%
Fatorando,
%
\begin{equation}
    (\omega_2-\omega_1)
    \left[
    -6t
    -
    \frac32
    (\omega_1^2+\omega_1\omega_2+\omega_2^2)
    \right]
    =0.
\end{equation}
%
Como $\omega_1\neq\omega_2$, pois estamos na curva de coexistência de fases,
%
\begin{equation}
    -6t
    -
    \frac32
    (\omega_1^2+\omega_1\omega_2+\omega_2^2)=0.
\end{equation}
%
Assim,
%
\begin{equation}    
\omega_1^2+\omega_1\omega_2+\omega_2^2=-4t.
\label{eq:pressure_condition}
\end{equation}
%
Agora, de acordo com a construção de Maxwell,
%
\begin{equation}
    \int_{\omega_1}^{\omega_2}
    \left[
    4t-6t\omega-\frac32\omega^3
    \right]\dd\omega
    =
    \pi(\omega_2-\omega_1).
\end{equation}
%
Integrando,
%
\begin{align}
    &4t(\omega_2-\omega_1)
    -
    3t(\omega_2^2-\omega_1^2)
    -
    \frac38(\omega_2^4-\omega_1^4)
    \nonumber\\
    &\hspace{2cm}
    =
    \pi(\omega_2-\omega_1).
\end{align}
%
Definamos
%
\begin{equation}
    S=\omega_1+\omega_2,
    \qquad
    P=\omega_1\omega_2.
\end{equation}
%
Então, $S^2 = (\omega_1+\omega_2)^2 = \omega_1^2 +2\omega_1\omega_2 +\omega_2^2 = \omega_1^2+ 2P +\omega_2^2$. Rearranjando essa expressão, segue que
%
\begin{equation}
    \omega_1^2+\omega_2^2
    =
    S^2-2P.
\end{equation}
%
Portanto,
%
\begin{align}
    \omega_1^2+\omega_1\omega_2+\omega_2^2
    &=
    (\omega_1^2+\omega_2^2)
    +\omega_1\omega_2
    \nonumber\\
    &=
    (S^2-2P)+P
    \nonumber\\
    &=
    S^2-P.
\end{align}
%
Então a relação \eqref{eq:pressure_condition} torna-se
%
\begin{equation}
    S^2-P=-4t,
\end{equation}
%
ou
%
\begin{equation}
    P=S^2+4t.
    \label{eq:P}
\end{equation}
%
%
Para simplificar a condição de Maxwell, podemos fatorar as diferenças
de potências que aparecem na expressão. Primeiramente,
%
\begin{equation}
    \omega_2^2-\omega_1^2
    =
    (\omega_2-\omega_1)(\omega_2+\omega_1)
    =
    (\omega_2-\omega_1)S.
\end{equation}
%
Da mesma forma, para a diferença das quartas potências, temos
%
\begin{align}
    \omega_2^4-\omega_1^4
    &=
    (\omega_2^2-\omega_1^2)
    (\omega_2^2+\omega_1^2)
    \nonumber\\
    &=
    (\omega_2-\omega_1)
    S
    (\omega_1^2+\omega_2^2).
\end{align}
%
Como já mostramos que
%
\begin{equation}
    \omega_1^2+\omega_2^2=S^2-2P,
\end{equation}
%
segue que
%
\begin{equation}
    \omega_2^4-\omega_1^4
    =
    (\omega_2-\omega_1)
    S(S^2-2P).
\end{equation}
%
Substituindo essas duas relações na condição de Maxwell, obtemos
%
\begin{align}
    &4t(\omega_2-\omega_1)
    -
    3t(\omega_2^2-\omega_1^2)
    -
    \frac38(\omega_2^4-\omega_1^4)
    =
    \pi(\omega_2-\omega_1)
    \nonumber\\
    &=
    4t(\omega_2-\omega_1)
    -
    3t(\omega_2-\omega_1)S
    -
    \frac38
    (\omega_2-\omega_1)
    S(S^2-2P)
    =
    \pi(\omega_2-\omega_1).
\end{align}
%
Para as duas fases coexistentes, temos $\omega_1\neq\omega_2$.
Portanto, podemos dividir toda a equação por
$(\omega_2-\omega_1)$:
%
\begin{equation}
    4t
    -3tS
    -\frac38S(S^2-2P)
    =
    \pi.
\label{eq:maxwell_SP}
\end{equation}
%
Podemos agora utilizar a relação da Eq.~\eqref{eq:P} na condição de Maxwell,
Eq.~\eqref{eq:maxwell_SP}. Para isso, calculamos
%
\begin{align}
    S^2-2P
    &=
    S^2-2(S^2+4t)
    \nonumber\\
    &=
    -S^2-8t.
\end{align}
%
Substituindo na Eq.~\eqref{eq:maxwell_SP}, obtemos
%
\begin{align}
    \pi
    &=
    4t-3tS
    -\frac38S(-S^2-8t)
    \nonumber\\
    &=
    4t-3tS
    +\frac38S^3
    +3tS.
\end{align}
%
Os termos lineares em $S$ contendo $t$ se cancelam, resultando em
%
\begin{equation}
    \pi
    =
    4t+\frac38S^3.
\label{eq:pi_S}
\end{equation}
%
%
A condição de Maxwell, juntamente com a igualdade das pressões,
permite obter uma segunda relação envolvendo $S$ e $P$. Como as
pressões nas duas fases são iguais,
%
\begin{equation}
    \pi(\omega_1)=\pi(\omega_2)=\pi,
\end{equation}
%
podemos tomar a média das duas expressões para a pressão:
%
\begin{align}
    \pi
    &=
    \frac{1}{2}
    \left[
    \pi(\omega_1)+\pi(\omega_2)
    \right]
    \nonumber\\
    &=
    \frac{1}{2}
    \left[
    4t-6t\omega_1-\frac32\omega_1^3
    +
    4t-6t\omega_2-\frac32\omega_2^3
    \right].
\end{align}
%
Usando $S=\omega_1+\omega_2$, obtemos
%
\begin{equation}
    \pi
    =
    4t-3tS
    -\frac34
    \left(
    \omega_1^3+\omega_2^3
    \right).
\end{equation}
%
A soma dos cubos pode ser escrita em termos de $S$ e $P$.
De fato,
%
\begin{align}
    \omega_1^3+\omega_2^3
    &= \omega_1^3 + 3\omega_1^2\omega_2 + 3\omega_1\omega_2^2 + \omega_2^3 -3\omega_1\omega_2(\omega_1+\omega_2) \\
    &=
    (\omega_1+\omega_2)^3
    -3\omega_1\omega_2(\omega_1+\omega_2)
    \nonumber\\
    &=
    S^3-3PS.
\end{align}
%
Consequentemente,
%
\begin{equation}
    \pi
    =
    4t-3tS
    -\frac34(S^3-3PS).
\end{equation}
%
Ou, equivalentemente,
%
\begin{equation}
    \pi
    =
    4t-3tS
    -\frac34S^3
    +\frac94PS.
\label{eq:pi_SP}
\end{equation}
%
Igualando as duas expressões Eqs.~\eqref{eq:pi_SP}
e~\eqref{eq:pi_S},temos
%
\begin{equation}
    4t-3tS
    -\frac34S^3
    +\frac94PS
    =
    4t+\frac38S^3.
\end{equation}
%
Cancelando o termo $4t$ dos dois lados,
%
\begin{equation}
    -3tS
    -\frac34S^3
    +\frac94PS
    =
    \frac38S^3.
\end{equation}
%
Multiplicando toda a equação por $8$, obtemos
%
\begin{equation}
    -24tS-6S^3+18PS=3S^3.
\end{equation}
%
Reorganizando os termos,
%
\begin{equation}
    -24tS-9S^3+18PS=0.
\end{equation}
%
Colocando $S$ em evidência,
%
\begin{equation}
    S\left(-24t-9S^2+18P\right)=0.
\end{equation}
%
Agora utilizamos a relação obtida anteriormente,
%
\begin{equation}
    P=S^2+4t.
\end{equation}
%
Substituindo $P$ na equação acima,
%
\begin{align}
    S\left[
    -24t-9S^2+18(S^2+4t)
    \right]
    &=0,
\end{align}
%
ou seja,
%
\begin{align}
    S\left[
    -24t-9S^2+18S^2+72t
    \right]
    &=0,
    \nonumber\\
    S\left[
    48t+9S^2
    \right]
    &=0.
\end{align}
%
Colocando um fator $3$ em evidência, finalmente obtemos
%
\begin{equation}
    3S(3S^2+16t)=0.
\end{equation}
%
Como o fator $3$ é diferente de zero, podemos escrever simplesmente
%
\begin{equation}
    S(3S^2+16t)=0.
\label{eq:S_condition}
\end{equation}
%
A Eq.~\eqref{eq:S_condition} possui duas raízes:
%
\begin{equation}
    S=0
    \qquad\text{ou}\qquad
    S^2=-\frac{16}{3}t.
\end{equation}
%
Estamos interessados na solução física que descreve as duas fases
coexistentes nas proximidades do ponto crítico. Nesse caso, a solução
simétrica é
%
\begin{equation}
    \boxed{S=0}.
\end{equation}
%
Como $S=\omega_1+\omega_2$, segue imediatamente que
%
\begin{equation}
    \omega_1+\omega_2=0,
\end{equation}
%
e, portanto,
%
\begin{equation}
    \omega_2=-\omega_1.
\end{equation}
%
Podemos então introduzir uma única variável $q$ para descrever os
deslocamentos das duas fases em relação ao ponto crítico:
%
\begin{equation}
    \omega_2=q,
    \qquad
    \omega_1=-q.
\end{equation}
%
Substituindo essas expressões na condição
\eqref{eq:pressure_condition}, temos
%
\begin{align}
    \omega_1^2+\omega_1\omega_2+\omega_2^2
    &=
    (-q)^2+(-q)(q)+q^2
    \nonumber\\
    &=
    q^2-q^2+q^2
    \nonumber\\
    &=
    q^2 = -4t.
\end{align}
%
Logo,
%
\begin{equation}
    q^2=-4t.
\end{equation}
%
Assim, para a coexistência das duas fases, é necessário que
$t<0$, isto é, que a temperatura esteja abaixo da temperatura crítica.
Escolhendo $q>0$, obtemos
%
\begin{equation}
    q=2\sqrt{-t}.
\end{equation}
%
Consequentemente, os volumes reduzidos das duas fases coexistentes
são dados por
%
\begin{equation}
    \boxed{
    \omega_1=-2\sqrt{-t},
    \qquad
    \omega_2=2\sqrt{-t}.
    }
\label{eq:omega_coexistence}
\end{equation}
%
Portanto, a separação entre os volumes das fases líquida e gasosa
próximas ao ponto crítico é
%
\begin{equation}
    \omega_2-\omega_1
    =
    4\sqrt{-t},
\end{equation}
%
mostrando explicitamente que a diferença entre os volumes das duas
fases desaparece como
%
\begin{equation}
    \boxed{
    \omega_2-\omega_1\propto(-t)^{1/2}
    }
\end{equation}
%
quando $t\rightarrow0^{-}$. Esse resultado fornece diretamente o
expoente crítico associado ao parâmetro de ordem da transição de fase
de van der Waals:
%
\begin{equation}
    \boxed{\beta=\frac12}.
\end{equation}
%
Substituindo
%
\begin{equation}
    \omega_1=-2\sqrt{-t}
\end{equation}
%
na equação de estado,
%
\begin{equation}
    \pi_{\mathrm{coex}}
    =
    4t
    -6t\omega_1
    -\frac32\omega_1^3.
\end{equation}
%
Como
%
\begin{equation}
    \omega_1^3 = \omega_1^2 \omega_1
    = 
    -4t\omega_1,
\end{equation}
%
temos
%
\begin{align}
    \pi_{\mathrm{coex}}
    &=
    4t
    -6t\omega_1
    +6t\omega_1
    \\
    &=
    4t.
\end{align}
%
Portanto, dentro da aproximação da expansão crítica considerada,
%
\begin{equation}
    \boxed{
    \pi_{\mathrm{coex}}=4t.
    }
    \label{eq:pcoex}
\end{equation}
%
Como
%
\begin{equation}
    p=p_c(1+\pi),
\end{equation}
%
a pressão ao longo da curva de coexistência é dada por
%
\begin{equation}
    \boxed{
    p_{\mathrm{coex}}
    =
    p_c(1+4t).
    }
\end{equation}
% ============================================================
\section*{Problema (6) -- Curva de coexistência e o locus de $v=v_c$: tangência no diagrama $p$--$T$}
%
Como obtido no exercício anterior, para $t<0$, a curva de coexistência nas proximidades do ponto crítico é
%
\begin{equation}
    \pi_{\mathrm{coex}}
    =
    4t.
\end{equation}
%
Portanto, a derivada da pressão reduzida ao longo da curva de coexistência é simplesmente
%
\begin{equation}
    \boxed{
    \left.
    \frac{\dd\pi_{\mathrm{coex}}}{\dd t}
    \right|_{t\to0^-}
    =4.
    }
\label{eq:derivative_coex}
\end{equation}
%
Para $t>0$, consideremos o locus
%
\begin{equation}
    v=v_c.
\end{equation}
%
Nesse caso, $\omega=0$. Na equação crítica (Eq.~\ref{eq:critical_expansion}), fazendo $\omega = 0$, fornece imediatamente
%
\begin{equation}
    \pi_{v=v_c}=4t.
\end{equation}
%
Sendo assim, utilizando a relação $p=p_c(1+\pi)$, 
%
\begin{equation}
    p_{v=v_c}
    =
    p_c(1+4t).
    \label{eq:locus}
\end{equation}
%
Por outro lado,
%
\begin{equation}
    \boxed{
    \left.
    \frac{\dd\pi_{v=v_c}}{\dd t}
    \right|_{t\to0^+}
    =4.
    }
\label{eq:derivative_vc}
\end{equation}
%
Portanto,
%
\begin{equation}
    \left.
    \frac{\dd\pi_{\mathrm{coex}}}{\dd t}
    \right|_{t\to0^-}
    =
    \left.
    \frac{\dd\pi_{v=v_c}}{\dd t}
    \right|_{t\to0^+}
    =4.
\end{equation}
%
Portanto, os coeficientes angulares das duas curvas coincidem no ponto crítico. Agora, lembramos que
%
\begin{equation}
    T=T_c(1+t),
    \qquad
    p=p_c(1+\pi).
\end{equation}
%
Diferenciando a primeira relação em relação a $t$, obtemos
%
\begin{equation}
    \frac{\dd T}{\dd t}=T_c,
\end{equation}
%
enquanto, da segunda,
%
\begin{equation}
    \frac{\dd p}{\dd t}
    =
    p_c\frac{\dd\pi}{\dd t}.
\end{equation}
%
Pela regra da cadeia,
%
\begin{equation}
    \frac{\dd p}{\dd T}
    =
    \frac{\dd p/\dd t}{\dd T/\dd t},
\end{equation}
%
e, portanto,
%
\begin{equation}
    \frac{\dd p}{\dd T}
    =
    \frac{p_c}{T_c}
    \frac{\dd\pi}{\dd t}.
\end{equation}
%
Para a curva de coexistência, usando
Eq.~\eqref{eq:derivative_coex}, temos
%
\begin{align}
    \left(\frac{\dd p}{\dd T}\right)_{\mathrm{coex}}
    &=
    \frac{p_c}{T_c}
    \left.
    \frac{\dd\pi_{\mathrm{coex}}}{\dd t}
    \right|_{t\to0^-}
    \nonumber\\
    &=
    \frac{p_c}{T_c}(4).
\end{align}
%
Assim,
%
\begin{equation}
    \boxed{
    \left(\frac{\dd p}{\dd T}\right)_{\mathrm{coex}}
    \longrightarrow
    \frac{4p_c}{T_c}
    \qquad
    (t\to0^-).
    }
\label{eq:slope_coex}
\end{equation}
%
Analogamente, ao longo da isócora crítica,
%
\begin{align}
    \left(\frac{\dd p}{\dd T}\right)_{v=v_c}
    &=
    \frac{p_c}{T_c}
    \left.
    \frac{\dd\pi_{v=v_c}}{\dd t}
    \right|_{t\to0^+}
    \nonumber\\
    &=
    \frac{p_c}{T_c}(4),
\end{align}
%
de modo que
%
\begin{equation}
    \boxed{
    \left(\frac{\dd p}{\dd T}\right)_{v=v_c}
    =
    \frac{4p_c}{T_c}.
    }
\label{eq:slope_vc}
\end{equation}
%
Finalmente, comparando as Eqs.~\eqref{eq:slope_coex}
e~\eqref{eq:slope_vc}, encontramos
%
\begin{equation}
    \boxed{
    \left(\frac{\dd p}{\dd T}\right)_{\mathrm{coex}}
    \longrightarrow
    \frac{4p_c}{T_c}
    =
    \left(\frac{\dd p}{\dd T}\right)_{v=v_c}.
    }
\end{equation}
%
Portanto, embora a curva de coexistência exista somente para
$T<T_c$, sua tangente aproxima-se continuamente da tangente da
isócora crítica quando $T\to T_c$. As duas curvas, portanto,
encontram-se \textbf{tangencialmente} no ponto crítico.
% ============================================================
\section*{Problema (7) -- Compressibilidade isotérmica de um fluido simples nas vizinhanças do ponto crítico}
%
A compressibilidade isotérmica de um fluido simples é definida por
%
\begin{equation}
    \kappa_T (T,v) 
    =
    -\frac{1}{v}
    \left(\pd{v}{p}\right)_T.
    \label{eq:kappa}
\end{equation}
%
Como a temperatura $T$ é mantida constante, podemos considerar
$v=v(p,T)$ e, localmente, $p=p(v,T)$ como funções inversas uma da
outra. Dessa forma, pela regra da função inversa para derivadas
parciais,
%
\begin{equation}
    \left(\pd{v}{p}\right)_T
    =
    \frac{1}{
    \left(\pd{p}{v}\right)_T
    }.
\end{equation}
%
Portanto,
%
\begin{equation}
    \kappa_T
    =
    -\frac{1}{v}
    \frac{1}{
    \left(\pdv{p}{v}\right)_T
    }.
\label{eq:kappa_inverse}
\end{equation}
%
Agora, vamos utilizar a equação de van der Waals para obter expressões assintóticas. O caminho $v=v_c$, $t\to0^+$, da expansão crítica, é dado por,
%
\begin{equation}
    \pi
    =
    4t
    -6t\omega
    -\frac32\omega^3+\mathcal{O}(\omega^4, t\omega^2)
\label{eq:critical_expansion_kappa}    
\end{equation}
%
Para calcular $\left(\pdv{\pi}{\omega}\right)_t$, devemos derivar
cada termo da Eq.~\eqref{eq:critical_expansion_kappa} em relação a
$\omega$, mantendo $t$ constante:
%
\begin{align}
    \left(\pdv{\pi}{\omega}\right)_t
    &=
    \pdv{}{\omega}
    \left[
    4t
    -6t\omega
    -\frac32\omega^3
    +\mathcal{O}(\omega^4, t\omega^2)
    \right]
    \nonumber\\
    &=
    -6t
    -\frac92\omega^2
    +\mathcal{O}(\omega^3).
\end{align}
%
Portanto, avaliando no ponto $\omega=0$ onde $v=v_c$, temos
%
\begin{equation}
    \left.
    \left(\pdv{\pi}{\omega}\right)_t
    \right|_{\omega=0}
    =
    -6t.
\label{eq:dpi_domega}
\end{equation}
%
Para relacionar essa derivada com
$\left(\pdv{p}{v}\right)_T$, utilizamos as definições das variáveis
reduzidas,
%
\begin{equation}
    p=p_c(1+\pi),
    \qquad
    v=v_c(1+\omega),
    \qquad
    T=T_c(1+t).
\end{equation}
%
Como $p_c$ e $v_c$ são constantes, a derivada de $p$ em relação a $v$,
mantendo $T$ constante, pode ser calculada diretamente pela regra da
cadeia:
%
\begin{equation}
    \left(\pdv{p}{v}\right)_T
    =
    \left(\pdv{p}{\pi}\right)_T
    \left(\pdv{\pi}{\omega}\right)_t
    \left(\pdv{\omega}{v}\right)_T.
\end{equation}
%
Cada um dos fatores pode ser calculado explicitamente. Das relações
$p=p_c(1+\pi)$ e $v=v_c(1+\omega)$, temos
%
\begin{equation}
    \left(\pdv{p}{\pi}\right)_T=p_c.
\end{equation}
%
e
%
\begin{equation}
    \dv{v}{\omega}=v_c,
\end{equation}
%
e, portanto,
%
\begin{equation}
    \left(\pdv{\omega}{v}\right)_T
    =
    \frac{1}{v_c}.
\end{equation}
%
Consequentemente,
%
\begin{align}
    \left(\pdv{p}{v}\right)_T
    &=
    p_c
    \left(\pdv{\pi}{\omega}\right)_t
    \frac{1}{v_c}
    \nonumber\\
    &=
    \frac{p_c}{v_c}
    \left(\pdv{\pi}{\omega}\right)_t.
\label{eq:p_and_pi_partial}    
\end{align}
%
Utilizando a Eq.~\eqref{eq:dpi_domega}, segue que
%
\begin{align}
    \left(\pdv{p}{v}\right)_{T,v=v_c}
    &=
    \frac{p_c}{v_c}
    \left.
    \left(\pdv{\pi}{\omega}\right)_t
    \right|_{\omega=0}
    \nonumber\\
    &=
    \frac{p_c}{v_c}(-6t).
\end{align}
%
Finalmente,
%
\begin{equation}
    \left(\pdv{p}{v}\right)_{T,v=v_c}
    =
    -\frac{6p_c}{v_c}t.
\label{eq:dpdv_critical}
\end{equation}
%
Consequentemente, substituindo esse resultado na expressão da compressibilidade isotérmica, obtemos
%
\begin{align}
    \kappa_T
    &=
    -\frac{1}{v_c}
    \frac{1}{
    (-\frac{6p_c}{v_c}t)
    }
    \nonumber\\
    &=
    \frac{1}{6p_c\,t}.
\end{align}
%
Portanto, a compressibilidade isotérmica diverge no ponto crítico como
%
\begin{equation}
    \boxed{
    \kappa_T(T,v_c)
    \sim
    \frac{1}{6p_c}t^{-1}.
    }
\end{equation}
%
Comparando com
%
\begin{equation}
    \kappa_T(T, v=v_c) \sim Ct^{-\gamma},
\end{equation}
%
nos permite identificar o expoente crítico e o prefator
%
\begin{equation}
    \boxed{\gamma=1},
    \qquad
    \boxed{C=\frac{1}{6p_c}}.
\end{equation}
%
Agora, para o caminho $T=T_c$, $p\to p_c^+$, implica que $t=0$.
%
Logo,
%
\begin{equation}
    \pi
    =
    -\frac32\omega^3
    +
    +\mathcal{O}(\omega^4, t\omega^2).
\end{equation}
%
Portanto,
%
\begin{equation}
    \omega
    =
    -\left(\frac{2\pi}{3}\right)^{1/3}.
\end{equation}
%
Além disso,
%
\begin{equation}
    \left.
    \left(\pdv{\pi}{\omega}\right)_t
    \right|_{t=0}
    =
    -\frac92\omega^2+\cdots.
\end{equation}
%
Assim, utilizando novamente as equações~\ref{eq:kappa_inverse} e \ref{eq:p_and_pi_partial},
%
\begin{align}
    \kappa_T(T=T_c, p)
    &=
    -\frac{1}{v_c}
    \frac{v_c}{p_c}
    \frac{1}{\left.
    \left(\pdv{\pi}{\omega}\right)_t
    \right|_{t=0}}
    \\
    &=
    -\frac{1}{p_c}
    \frac{1}{(-\frac92\omega^2)}
    \\
    &=
    \frac{2}{9p_c\omega^2}.
\end{align}
%
Como
%
\begin{equation}
    \omega^2
    =
    \left(\frac{2\pi}{3}\right)^{2/3},
\end{equation}
%
temos
%
\begin{align}
    \kappa_T(T=T_c, p)
    &=
    \frac{2}{9p_c}
    \left(\frac{3}{2}\right)^{2/3}
    \pi^{-2/3}.
\end{align}
%
Logo,
%
\begin{equation}
    \boxed{
    \kappa_T(T_c,p)
    \sim
    \bar C
    \left(
    \frac{p-p_c}{p_c}
    \right)^{-\gamma_p},
    }
\end{equation}
%
com
%
\begin{equation}
    \boxed{\gamma_p=\frac23}
\end{equation}
%
e
%
\begin{equation}
    \boxed{
    \bar C
    =
    \frac{2}{9p_c}
    \left(\frac32\right)^{2/3}.
    }
\end{equation}
%
quando $p\to p_c^+$. Portanto, os valores obtidos para os expoentes críticos são
%
\begin{equation}
    \boxed{
    \gamma=1,
    \qquad
    \gamma_p=\frac23.
    }
\end{equation}
%
O valor do expoente depende do caminho de aproximação ao ponto crítico, enquanto os valores dos expoentes críticos é universal dentro da classe de universalidade de campo médio de van der Waals.
%%%%%%%%%%%%%%%%%%%%%%%%%%%%%%%%%%%%%%%%%%%%%%%%%%%%%%%%%%%%%%%%%%%%%%%%%%%%%%%%
\section*{Problema (8) -- Forma assintótica do calor específico a volume constante}
%%%%%%%%%%%%%%%%%%%%%%%%%%%%%%%%%%%%%%%%%%%%%%%%%%%%%%%%%%%%%%%%%%%%%%%%%%%%%%%%
Neste problema queremos determinar o comportamento assintótico do calor
específico a volume constante nas proximidades do ponto crítico do fluido de
van der Waals, considerando o caminho
%
\begin{equation}
    v=v_c.
\end{equation}
Em particular, devemos analisar separadamente os casos $t>0$ e $t<0$,
mostrar que o calor específico permanece finito no ponto crítico e determinar
a possível descontinuidade entre os limites $t\to0^+$ e $t\to0^-$.
%
Definimos a temperatura reduzida por
\begin{equation}
    t\equiv\frac{T-T_c}{T_c}.
    \label{eq:t_reduced_cv}
\end{equation}
%
Assim,
%
\begin{equation}
    T=T_c(1+t),
    \qquad
    \frac{\dd t}{\dd T}=\frac{1}{T_c}.
    \label{eq:t_derivative_cv}
\end{equation}
%
Para o fluido de van der Waals, a energia livre específica pode ser escrita
na forma
%
\begin{equation}
    f_{\rm vdW}(T,v)
    =
    -RT\ln(v-b)-\frac{a}{v}+f_0(T),
    \label{eq:vdw_free_energy_cv}
\end{equation}
%
onde $f_0(T)$ é uma função exclusivamente da temperatura. Essa função não é determinada pela equação de estado de van der Waals e terá um papel importante na discussão do calor específico.
%
A entropia específica é obtida a partir da energia livre de Helmholtz através da relação termodinâmica
%
\begin{equation}
    s
    =
    -\left(\pd{f}{T}\right)_v.
    \label{eq:entropy_free_energy}
\end{equation}
%
Aplicando essa relação à energia livre de van der Waals, obtemos
%
\begin{align}
    s(T,v)
    &=
    -\left[
    -R\ln(v-b)+f_0'(T)
    \right]
    \nonumber\\
    &=
    R\ln(v-b)-f_0'(T).
    \label{eq:entropy_vdw}
\end{align}
%
O calor específico a volume constante é definido por
%
\begin{equation}
    c_V
    =
    T\left(\pd{s}{T}\right)_v.
    \label{eq:cv_definition}
\end{equation}
%
Portanto,
%
\begin{align}
    c_V
    &=
    T
    \left[
    -f_0''(T)
    \right]
    \nonumber\\
    &=
    \boxed{
    -T f_0''(T).
    }
    \label{eq:cv_f0_general}
\end{align}
%
Para a análise crítica, escreveremos
%
\begin{equation}
    c_V^+
    \equiv
    \lim_{t\to0^+}c_V,
\end{equation}
%
representando o limite do calor específico quando o ponto crítico é
aproximado pelo lado de temperaturas superiores.
%
Para estudar o comportamento crítico, introduzimos também a variável reduzida de volume
%
\begin{equation}
    \omega
    \equiv
    \frac{v-v_c}{v_c},
    \label{eq:omega_cv}
\end{equation}
%
e a pressão reduzida
%
\begin{equation}
    \pi
    \equiv
    \frac{p-p_c}{p_c}.
    \label{eq:pi_cv}
\end{equation}
%
Nas proximidades do ponto crítico, a equação de estado de van der Waals pode ser expandida na forma
%
\begin{equation}
    \pi
    =
    4t-6t\omega-\frac{3}{2}\omega^3
    +\mathcal{O}(t\omega^2,\omega^4).
    \label{eq:critical_eos_cv}
\end{equation}
% 
(ii) Consideremos inicialmente o caminho    $t>0$, $v=v_c$, onde $\omega=0$. Integrando em relação a $\omega$, obtemos uma função proporcional a energia livre, 
%
\begin{equation}
    f(t,\omega)
    =
    4t\omega
    -3t\omega^2
    -\frac38\omega^4
    + f_0(t)
    +\cdots.
    \label{eq:free_energy_critical}
\end{equation}
%
Essa parte da energia livre se anula para o caminho considerado, somente restando a função arbitrária. 
%
\begin{equation}
    f(t,\omega=0)=f_0(t).
\end{equation}
%
Assim, para a equação de estado de van der Waals, o calor específico
a volume constante permanece finito no ponto crítico. Em termos do
expoente crítico $\alpha$, escrevemos
%
\begin{equation}
    c_V(T,v=v_c)
    \sim
    t^{-\alpha},
    \qquad
    t\to0^+,
\end{equation}
%
com
%
\begin{equation}
    \boxed{\alpha=0}.
    \label{eq:cv_plus}
\end{equation}
%
(i) Agora, para determinar o comportamento abaixo de $T_c$ ($t < 0$), precisamos fazer uma construção de Maxwell para eliminar as alças de van der Waals restaurando a convexidade da energia livre. Sendo assim, temos que
%
\begin{equation}
    \int_{v_l}^{v_g}
    \left[
    p(T,v)-p_{\mathrm{coex}}(T)
    \right]\dd v
    =
    0,
\end{equation}
%
onde $l$ e $g$ denotam, respectivamente, as fases líquida e gasosa na curva de coexistência. A condição de Maxwell garante que as duas fases possuem a mesma energia livre na coexistência. A energia livre de Helmholtz $f=f(T,v)$ deve obedecer à relação
termodinâmica
%
\begin{equation}
    p
    =
    -\left(\pd{f}{v}\right)_T.
    \label{eq:p_free_energy_cv}
\end{equation}
%
%
Agora, abaixo da temperatura crítica, na região de coexistência, a energia livre corrigida é dada pela construção da tangente comum (obtida no problema (3)),
%
\begin{equation}
    f_{\rm corr}(T,v)
    =
    f_{\rm vdW}(T,v_l)
    +
    \frac{v-v_l}{v_g-v_l}
    \left[
    f_{\rm vdW}(T,v_g)
    -
    f_{\rm vdW}(T,v_l)
    \right].
    \label{eq:maxwell_free_energy_cv}
\end{equation}
%
em que $f_{\rm vdW}(T,v_l)$ e $f_{\rm vdW}(T,v_g)$ são as energias livres de van der Waals sem a construção de Maxwell, avaliadas nos volumes de coexistência. Então, fixando o volume $v = v_c$, temos
%
\begin{equation}
    f_{\rm corr}(T,v=v_c)
    =
    f_{\rm vdW}(T,v_l)
    +
    \frac{v_c-v_l}{v_g-v_l}
    \left[
    f_{\rm vdW}(T,v_g)
    -
    f_{\rm vdW}(T,v_l)
    \right].
    \label{eq:maxwell_free_energy_cv_critical}
\end{equation}
%
Nas proximidades do ponto crítico, os volumes de coexistência são (como obtido no problema (5))
%
\begin{equation}
    \omega_l=-2\sqrt{-t},
    \qquad
    \omega_g=2\sqrt{-t},
    \qquad
    t<0,
\end{equation}
%
de modo que
%
\begin{equation}
    v_l=v_c(1+\omega_l),
    \qquad
    v_g=v_c(1+\omega_g).
\end{equation}
%
Consequentemente,
%
\begin{align}
    v_c-v_l
    &=
    -v_c\omega_l
    =
    2v_c\sqrt{-t},
    \\
    v_g-v_l
    &=
    v_c(\omega_g-\omega_l)
    =
    4v_c\sqrt{-t}.
\end{align}
%
Portanto,
%
\begin{equation}
    \frac{v_c-v_l}{v_g-v_l}
    =
    \frac12.
\end{equation}
%
A Eq.~\eqref{eq:maxwell_free_energy_cv_critical} reduz-se, portanto,
a
%
\begin{equation}
    f_{\rm corr}(T,v_c)
    =
    \frac12
    \left[
    f_{\rm vdW}(T,v_l)
    +
    f_{\rm vdW}(T,v_g)
    \right].
    \label{eq:fcorr_average_cv}
\end{equation}
%
Para obter o comportamento próximo ao ponto crítico, utilizamos a
expansão da energia livre correspondente à equação de estado
%
\begin{equation}
    f_{\rm vdW}(t,\omega)
    =
    f_0(t)
    +
    4t\omega
    -3t\omega^2
    -\frac38\omega^4
    +\cdots.
    \label{eq:fvdw_expansion_cv}
\end{equation}
%
Substituindo $\omega_l=-2\sqrt{-t}$, obtemos
%
\begin{align}
    f_{\rm vdW}(t,\omega_l)
    &=
    f_0(t)
    +4t(-2\sqrt{-t})
    -3t(-2\sqrt{-t})^2
    -\frac38(-2\sqrt{-t})^4
    +\cdots
    \nonumber\\
    &=
    f_0(t)
    -8t\sqrt{-t}
    +6t^2
    +\cdots.
\end{align}
%
Analogamente, para $\omega_g=2\sqrt{-t}$,
%
\begin{align}
    f_{\rm vdW}(t,\omega_g)
    &=
    f_0(t)
    +4t(2\sqrt{-t})
    -3t(2\sqrt{-t})^2
    -\frac38(2\sqrt{-t})^4
    +\cdots
    \nonumber\\
    &=
    f_0(t)
    +8t\sqrt{-t}
    +6t^2
    +\cdots.
\end{align}
%
Substituindo essas expressões na Eq.~\eqref{eq:fcorr_average_cv},
%
\begin{align}
    f_{\rm corr}(t,v_c)
    &=
    \frac12
    \left[
    f_0(t)
    -8t\sqrt{-t}
    +6t^2
    +
    f_0(t)
    +8t\sqrt{-t}
    +6t^2
    \right]
    +\cdots
    \nonumber\\
    &=
    f_0(t)+6t^2+\cdots.
\end{align}
%
Portanto,
%
\begin{equation}
    f_{\rm corr}(t,v_c)
    =
    f_0(t)+6t^2+\mathcal{O}(|t|^{5/2}),
    \qquad t\to0^-.
    \label{eq:fcorr_negative_cv}
\end{equation}
%
Observe que os termos não analíticos proporcionais a
$t\sqrt{-t}$ das fases líquida e gasosa possuem sinais opostos e, portanto, cancelam-se na construção de Maxwell.
%
A entropia molar é dada por
%
\begin{equation}
    s
    =
    -\left(\pd{f}{T}\right)_v.
\end{equation}
%
Como
%
\begin{equation}
    t=\frac{T-T_c}{T_c},
    \qquad
    \pdv{t}{T}=\frac{1}{T_c},
\end{equation}
%
temos
%
\begin{equation}
    \pdv{}{T}
    =
    \frac{1}{T_c}\pdv{}{t}.
\end{equation}
%
Assim, para $t<0$,
%
\begin{align}
    s_-(t)
    &=
    -\frac{1}{T_c}
    \pdv{}{t}
    \left[
    f_0(t)+6t^2+\cdots
    \right]
    \nonumber\\
    &=
    -\frac{1}{T_c}
    \left[
    f_0'(t)+12t+\cdots
    \right].
\end{align}
%
Portanto,
%
\begin{equation}
    \boxed{
    s_-(t)
    =
    -\frac{f_0'(t)}{T_c}
    -\frac{12}{T_c}t+\cdots.
    }
\end{equation}
%
Por outro lado, para $t>0$, ao longo do caminho $v=v_c$, temos
%
\begin{equation}
    f(t,v_c)=f_0(t),
\end{equation}
%
e, portanto,
%
\begin{equation}
    \boxed{
    s_+(t)
    =
    -\frac{f_0'(t)}{T_c}.
    }
\end{equation}
%
Consequentemente, a entropia permanece contínua no ponto crítico:
%
\begin{equation}
    \lim_{t\to0^-}s(t)
    =
    \lim_{t\to0^+}s(t)
    =
    -\frac{f_0'(0)}{T_c}.
\end{equation}
%
O calor específico a volume constante é
%
\begin{equation}
    c_V
    =
    -T
    \left(\pdv[2]{f}{T}\right)_v
    =
    -\frac{T}{T_c^2}
    \left(\pdv[2]{f}{t}\right)_v.
\end{equation}
%
Para $t<0$, segue da Eq.~\eqref{eq:fcorr_negative_cv} que
%
\begin{equation}
    \left(\pdv[2]{f_{\rm corr}}{t}\right)_{v_c}
    =
    f_0''(t)+12+\cdots,
\end{equation}
%
de modo que
%
\begin{equation}
    c_V^{(-)}
    =
    -\frac{T}{T_c^2}
    \left[
    f_0''(t)+12+\cdots
    \right].
\end{equation}
%
No limite $t\to0^-$,
%
\begin{equation}
    \boxed{
    c_V^{(-)}
    \longrightarrow
    -\frac{1}{T_c}
    \left[
    f_0''(0)+12
    \right].
    }
    \label{eq:cv_minus}
\end{equation}
%
Enquanto, para $t\to0^+$,
%
\begin{equation}
    \boxed{
    c_V^{(+)}
    \longrightarrow
    -\frac{1}{T_c}f_0''(0).
    }
    \label{eq:cv_plus_final}
\end{equation}
%
Assim, o calor específico permanece finito no ponto crítico, mas
apresenta uma descontinuidade finita:
%
\begin{equation}
    \boxed{
    \Delta c_V
    \equiv
    c_V^{(-)}-c_V^{(+)}
    =
    -\frac{12}{T_c}.
    }
\end{equation}
%
Portanto, não há divergência do calor específico no ponto crítico.
Seu comportamento assintótico é
%
\begin{equation}
    \boxed{
    c_V\sim |t|^{-\alpha},
    \qquad
    \alpha=0.
    }
\end{equation}
%
%Agora, fazendo uma expansão em termos da temperatura reduzida t e notando que os volumes $v_c, v_l, v_g$ estão fixos e que $f_{\rm corr}(T,v=v_c)$ depende apenas da temperatura, obtemos 
%

%
%Ainda, fazendo expansões para $t \rightarrow0$ das energias livres de van der Waals, $f_{\rm vdW}(T,v_g)$ e $f_{\rm vdW}(T,v_l)$, chegamos em, 
%
%
%Assim, os resultados assintóticos para a entropia e o calor específico são dados por,
%
%
%%%%%%%%%%%%%%%%%%%%%%%%%%%%%%%%%%%%%%%%%%%%%%%%%%%%%%%%%%
\begin{thebibliography}{99}
\bibitem{salinas} S. R. A. Salinas, \emph{Introdução à Física Estatística}, 2ª ed., 3ª reimpressão, Edusp, São Paulo, 2013.
%\bibitem{gibbs} J. W. Gibbs, \emph{Elementary Principles in Statistical Mechanics}, Charles Scribner's Sons, 1902.
%\bibitem{stanley} H. E. Stanley, \emph{Introduction to Phase Transitions and Critical Phenomena}, Oxford University Press.
%\bibitem{yeomans} J. M. Yeomans, \emph{Statistical Mechanics of Phase Transitions}, Oxford University Press.
%\bibitem{pathria} R. K. Pathria and P. D. Beale, \emph{Statistical Mechanics}, Academic Press.
%\bibitem{huang} K. Huang, \emph{Statistical Mechanics}, Wiley.
%\bibitem{goldenfeld} N. Goldenfeld, \emph{Lectures on Phase Transitions and the Renormalization Group}, CRC Press.
\end{thebibliography}
%%%%%%%%%%%%%%%%%%%%%%%%%%%%%%%%%%%%%%%%%%%%%%%%%%%%%%%%%%
\end{document}
